% ============================================================
%  Georg Brandes, Samlede Skrifter. Trettende Bind (Supplementbind).
%  Kjøbenhavn: Gyldendalske Boghandels Forlag (F. Hegel & Søn).
%  Græbes Bogtrykkeri.  1903.
%  Review of H. Brøchner, Bidrag til Opfattelsen af Filosofiens historiske Udvikling (1869).
%  TRANSCRIPTION (Danish), printed pp. 115--123.
%
%  One of three transcription directories cut from this volume:
%      striden/               pp. 1--42    Indledning; Striden om Tro og Viden
%      nielsen-anmeldelser/   pp. 85--109  five pieces on Nielsen, 1867
%      brochner-bidrag/       pp. 115--123 review of Brøchner's Bidrag (1869)
%  NOT transcribed: pp. 43--84 (Dualismen i vor nyeste Filosofi), 110--114,
%  124 ff.  See the note on pp. 43--84 below.
%
%  ---- DATE -------------------------------------------------------------
%  The title page (PDF 7) reads 1903.  The Indledning (p. 3) is signed
%  "April 1903." / "G. B."  (catalog.yaml's 1899 is the series date, not
%  this volume's.)  This is Brandes's 1903 revision of pieces first printed
%  1865--69, in the orthography of 1903 (Filosofi, Teologi, eksakt, hv-,
%  single vowels), transcribed as printed.  Nothing has been harmonised
%  with the first printings or with the 1866 Dualismen.
%
%  ---- SOURCE -----------------------------------------------------------
%  bibliotek/Brandes, Georg/1903-samlede-skrifter-13.pdf
%      Google Books (University of Chicago copy), 552 PDF pp., 600 dpi
%      bitonal JBIG2, embedded text layer (witness 1).
%      sha256 9579f1fd8672884bcc7e149b43f5c8c640dcaaefa1aceca01d5f89cbed016704
%  Witness 2 (never copy-text): Project Runeberg's OCR text of the same
%      edition (gbsamskr/13), bibliotek/Brandes, Georg/gbsamskr-13.txt,
%      cut into printed pages by alignment against witness 1 (per-page word
%      agreement between the two witnesses, pp. 115--123: mean 0.963, min 0.937 over 9 pages).
%      It is a different OCR engine's reading; whether Runeberg photographed
%      a different physical copy is not known.
%
%  ---- PAGE MAP ---------------------------------------------------------
%      printed p. = PDF p. - 12, uniform over printed 1--136.
%  Verified 2026-09-22 AT THE IMAGE: folio read off the rendered page at
%  printed 9, 29, 38, 42, 44, 86, 97, 104, 109, 116, 123 (all match; pp. 2
%  and 17 print no folio and were placed by content).  Independently read
%  by H. Halvorson at 8, 9, 42, 43, 54, 85, 93, 123, 136.
%
%  ---- LAYOUT OF THE VOLUME ---------------------------------------------
%  Every piece begins on a fresh page with a sunk head: title in large
%  italic, the year of first publication in roman in parentheses, a short
%  rule; the opening page carries no folio and no running head.  Every piece
%  ends with a short centred rule (\piecerule).  Running heads carry a short
%  title of the piece (a third reading, recorded at each piece's head where
%  it differs); they are not transcribed.
%
%  ---- CONVENTIONS ------------------------------------------------------
%  Diplomatic.  Orthography, punctuation and capitalisation as printed.
%  Emphasis: ITALIC is the volume's only emphasis device -- no letterspacing
%    and no bold were found in pp. 1--123 -- and every italic span is set
%    \emph{}, whether it marks emphasis, a title or a foreign word; no
%    attempt is made to distinguish the functions.
%  Quotation marks: the volume quotes with SINGLE guillemets pointing
%    outward, ‹...› (U+2039/U+203A), for quotations, titles and terms alike;
%    double «...» occur only in a head (p. 102).  Transcribed as printed.
%  Dashes: spaced em-dash, set ---.  ɔ: as printed.
%  Footnotes: *), set with \fnmark{*)}{...} at the mark.
%  A word broken across a printed line is joined (Cau\-%); a hyphen printed
%    at a line end as part of the word is kept.
%  Printer's errors are transcribed AS PRINTED and logged in a % comment at
%    the site; nothing is silently corrected.
%  Punctuation is as the Chicago copy prints it.  The 600 dpi bitonal image
%    sometimes loses the tail of a faint comma; where it shows a point but
%    witness 1 (Google's OCR, made from its capture BEFORE compression to
%    bitonal JBIG2) reads a comma, the comma is transcribed and commented;
%    where the image and witness 1 both show a point, the point is
%    transcribed and logged as a printer's error, whatever the syntax wants.
%  Letters the bitonal scan may not resolve (ø/o, æ/œ, worn e/c): where the
%    printed form would not be a word, the word is read unless the glyph
%    positively shows otherwise (TRANSCRIPTION-PLAYBOOK §2 rule 3); each case
%    is commented.  In this fount the italic æ has a closed o-like bowl and no
%    separate a-stem, so it looks like œ (see p. 27 `duæ', p. 10 Fædrelandet).
%  A hyphen at a line end that could be either a compound hyphen or a word
%    break is resolved by the volume's own usage MID-LINE: printed only solid
%    -> joined; only hyphenated -> kept; both or never -> the printed hyphen
%    is kept.  Each case is commented with its evidence.
%  A word divided across a PAGE turn is joined with `%' before the page
%    marker (e.g. Spørgs%  / \opage{122}maal); the divisional hyphen is not set.
%
%  ---- CONTENTS vs HEADS ------------------------------------------------
%  The Indhold (PDF 9) gives short titles; the heads over the pieces name
%  the book under review and the year.  Both readings are kept: the Indhold
%  reading is the first argument of \piecehead (table of contents, running
%  head), the head as printed is the second.  Recorded at each piece.
%
%  ---- READINGS OF THE HEADS ----------------------------------------------
%  Indhold (PDF 9): `Filosofiens historiske Udvikling 115'.
%  Head as printed: `H. Brøchner: Bidrag til Opfattelsen af Filosofiens' /
%  `historiske Udvikling' / `(1869)'.  The running head on p. 122 prints
%  `Filosoflens' (running heads are not transcribed).
%
%  ---- PRINTER'S ERRORS AND DOUBTFUL READINGS (each commented at its site) --
%  p. 116 `kans Talent'.  p. 117 `Livsanskuelse', `urimelig': worn e that
%  prints like c (sense reading); `anti-/rationelle' joined (the volume
%  prints `antirationel' solid mid-line).  p. 118 `dens Ophavs, til' as
%  printed.  p. 120 `Farve, Saaledes', `Folkeaand, Naar' (comma before a
%  new sentence).  p. 121 `fremstilies'.
%  Guillemet balance 0.
%
%  ---- pp. 43--84, NOT TRANSCRIBED --------------------------------------
%  pp. 43--84 reprint "Dualismen i vor nyeste Filosofi" (1866).  The 1866
%  first edition is already in this collection (texts/brandes/dualismen),
%  fully transcribed and collated against two scans.  The 1903 text is
%  Brandes's own revision, in modernised orthography and with altered
%  chapter titles (Indhold: I. Indledning 43, Den absolute Uensartethed 54,
%  Det Antirationelle 67, Dualismens Følger 71), so it is a VARIANT WITNESS
%  to the 1866 text and deserves a collation of its own, not a second
%  transcription.
%
%  ---- HOW THIS TEXT'S ACCURACY WAS ESTABLISHED ------------------------
%  Word-accuracy: the batch was diffed against TWO machine witnesses before
%  splicing (ocrdiff.py; TRANSCRIPTION-PLAYBOOK §3 step 4) -- witness 1 the
%  embedded layer, witness 2 Runeberg's OCR -- and then checked with
%  consensus.py, which compares case-sensitively and with punctuation where
%  the two witnesses agree with each other.  1 batch; witness 1: 1
%  candidate, witness 2: 5, in both: 0, consensus: 0; 0 transcription errors
%  surfaced; 0 unresolved.
%  Independent collation: a seeded random sample of 6 of the 9 pages
%  (seed 20260922: pp. 116, 119; top-up seed 20260922-top-up: pp. 115, 118,
%  120, 121) was read line by line at the image by readers who had not
%  transcribed them.  0 discrepancies.  The review of the batch agent's
%  doubtful calls overturned one (p. 117 `anti-/rationelle', kept -> joined,
%  by the rule for line-end hyphens); one italic call spot-checked, confirmed.
% ============================================================

\documentclass[12pt,a4paper,openany]{book}
\usepackage[utf8]{inputenc}
\usepackage[T1]{fontenc}
\usepackage{libertinus}
\usepackage{libertinust1math}
\usepackage{graphicx}    % for \reflectbox (the old Danish ɔ: id-est mark)
\DeclareUnicodeCharacter{0254}{\reflectbox{c}}  % ɔ (U+0254) = reversed c
\usepackage[danish]{babel}
\usepackage[protrusion=true,expansion=true]{microtype}
\usepackage{setspace}
\usepackage[margin=1.2in]{geometry}
\usepackage{fancyhdr}
\setlength{\headheight}{28pt}

% Footnotes as the 1903 printing sets them: *), not 1, 2, 3.  Set the mark
% explicitly at each site with \fnmark{*)}{...} (or \fnmark{**)}{...} should a
% page carry a second note); the default is restored afterwards.
\renewcommand{\thefootnote}{*)}
\newcommand{\fnmark}[2]{\renewcommand{\thefootnote}{#1}\footnote{#2}%
  \renewcommand{\thefootnote}{*)}}

% A piece's head.  #1 = the title as the volume's INDHOLD prints it (used for
% the table of contents and the running head); #2 = the head block as printed
% over the piece, with its own \opage.  The two readings differ, and both are
% kept (TRANSCRIPTION-PLAYBOOK §6).  The short rule under every head is
% printed, and is set by the macro.
\newcommand{\piecehead}[2]{\clearpage\phantomsection
  \addcontentsline{toc}{section}{#1}\markboth{#1}{#1}%
  \begin{center}\large #2\\[6pt]\rule{2cm}{0.4pt}\end{center}\par\medskip}
% The short centred rule printed at the end of every piece.
\newcommand{\piecerule}{\par\begin{center}\rule{1.5cm}{0.4pt}\end{center}\par}

\usepackage{hyperref}
\hypersetup{pdftitle={Brandes, Samlede Skrifter XIII: Filosofiens historiske Udvikling (1869)}, pdfauthor={Georg Brandes}, hidelinks}

\pagestyle{fancy}
\fancyhf{}
\fancyhead[LE,RO]{\thepage}
\fancyhead[RE]{\textit{Samlede Skrifter XIII (1903)}}
\fancyhead[LO]{\textit{\leftmark}}
\setstretch{1.3}

\title{\textbf{Filosofiens historiske Udvikling (1869)}\\[8pt]\large Georg Brandes, \textit{Samlede Skrifter}, Trettende Bind (Supplementbind),\\ pp.~115--123}
\author{}
\date{Kjøbenhavn: Gyldendalske Boghandels Forlag (F.\ Hegel \& Søn), 1903}

\usepackage{marginnote}
\newcommand{\opage}[1]{\marginnote{\footnotesize\textit{[#1]}}}
\newcommand{\apage}[1]{\marginnote{\footnotesize\textit{[#1]}}}
\begin{document}
\maketitle
\thispagestyle{empty}
\tableofcontents

\mainmatter

% ============================================================
%  Filosofiens historiske Udvikling (printed pp. 115--123)
%  Indhold (PDF 9) reads: `Filosofiens historiske Udvikling 115'
%  Head as printed: `H. Brøchner: Bidrag til Opfattelsen af Filosofiens'
%    / `historiske Udvikling' (both lines italic) / `(1869)' (roman),
%    short rule.  p. 115 has no folio and no running head; its foot
%    carries the signature `8*' -- not transcribed.
%  Running head (pp. 116--123) reads: `Filosofiens historiske Udvikling',
%    the Indhold reading; p. 122 prints it `Filosoflens' (broken fi
%    ligature in the head only -- running heads are not transcribed).
%  No dateline or signature; the piece ends on p. 123 with a short rule.
%  Printer's defects, transcribed as printed and logged at the site:
%    p. 116 `kans Talent' (for hans); p. 120 `Farve, Saaledes' and
%    `Folkeaand, Naar' (comma before a new sentence, twice);
%    p. 121 `fremstilies' (for fremstilles).
% ============================================================
% --- p. 115 ---
\piecehead{Filosofiens historiske Udvikling}{\opage{115}\emph{H. Brøchner: Bidrag til Opfattelsen af Filosofiens}\\
\emph{historiske Udvikling}\\[2pt](1869)}
Man læser i Almindelighed med en vis Forundring, at de
betydeligste Mænd i Fortiden som oftest netop har været de,
paa hvilke man i deres Samtid mindst har givet Agt; man
undrer sig ved at erfare, hvilke Mænd der har været foretrukne
for dem, hvilken Tidsstrømning de har havt imod sig, hvor
ensom deres Stilling har været, og man føler sig forvisset om,
at i vore Dage er Sligt forbi, at vor Tid forstaar at finde og
gribe Fortjenesten i dens skjulteste Vraa og drage den frem for
Lyset; man smigrer sig med den Tanke, at de Navne, som i
vore Dage lyder hyppigst og højest, ogsaa vil være de, som
Eftertiden har paa Læben. Ikke mange Ting i Verden er saa
visse, som at man heri tager sørgeligt fejl. Navne, som nu fej\-%
res, Personer, som nu Hvermand kender, vil være Eftertiden et
Intet, og Personligheder, som nu kendes og skattes af en lille
Kres, vil efterlade sig blivende Spor af deres Liv og Arbejde.
Men værre er Den dog faren, som overlever sin Hæder, end
den, hvis Hæder vokser saa langsomt, at den først, naar Tiden
er kommen i Højde med Personligheden, staar med tydelige
Omrids.

Professor \emph{Brøchner} er en Mand, som hører til den sidste
Klasse; han er ikke nogen almenyndet Mand; han har aldrig
gjort noget Skridt mod Populariteten, og han er altfor langt
forud for sine Omgivelser til at skattes af dem; ti man skatter
% --- p. 116 ---
\opage{116}kun, hvad man forstaar. Han har virket i Stilhed, og var han
bleven bortreven af Sygdom, før han endelig i Fjor begyndte
Udgivelsen af sine Skrifter, da vilde neppe en halv Snes Men\-%
nesker have vidst, at Landet havde mistet sin grundigste og
skarpeste Tænker. Lykkeligvis har han overstaaet sin lange
Sygelighed, og den af hans faa Disciple med saa stor Længsel
imødesete Udgivelse af hans Arbejder er paabegyndt.

Naar man betragter den Bog, hvis Titel ovenfor er nævnt,
da er Ens første Indtryk en smertelig Følelse. Er det ikke
sørgeligt, at Danmark efterhaanden er indskrænket til en saa\-%
dan Lidenhed, at de bedste Skrifter ingen Læsere kan finde,
ingen Anerkendelse kan høste, intet Udbytte kan give og ingen
Kritik kan modtage? Var det ikke bedre baade for Videnskabs\-%
manden selv og for hele den europæiske Dannelse, om Mænd
som Professor Brøchner og vore andre betydelige Videnskabs\-%
mænd levede i et af de store Hovedlande, udvikledes gennem
Kritik og modtog Tilskyndelser fra alle Sider? Maaske var det
bedre. Men det forekommer mig, at en saadan Mands Lod og\-%
saa under de nuværende Omstændigheder har sin Skønhed.
Han og hans Lige er den store europæiske Kulturs yderst
fremskudte Vedetter eller Forposter i det høje Norden; deres
Stilling er mest udsat; ti Forposttjenesten er streng og farlig;
men deres videnskabelige Liv udkræver netop derved en Ka\-%
rakterfasthed, et Mod, til visse Tider en Soldater-Dristighed,
som mindre ualmindelige Forhold ikke udvikler, og som indgyder
% `kans' for `hans': printer's error, transcribed as printed (both
% witnesses and the image, 600 dpi, agree on k).
Kærlighed til Manden, smykker og styrker kans Talent.

Professor Brøchners Lærervirksomhed ved Universitetet er,
skønt flereaarig, endnu ikke saa lang, at dens Indflydelse tyde\-%
ligt kan spores; men at den er betydelig, det er vist, og derom vil
Enhver, der kender mange yngre Studerende, kunne vidne. Pro\-%
fessor Brøchner er den første for al Teologi befriede Filosof,
som vort Universitet har besiddet; naar G. Sibbern, der har
maattet bøde haardt for sin Frisindethed, undtages, er han til\-%
lige den eneste nulevende danske. Han blev ansat paa en Tid,
da de teologiske Indflydelser af den forskelligste Art var mæg\-%
tige herhjemme; han har talt fra sit ensomme Kateder for sit
ikke altid talrige Auditorium, medens rundt omkring de Alt
overdøvende, aandelig reaktionære Røster raabte mellem hinan\-%
den og forvirrede enhver tænkende Hjerne. Han har holdt
fast ved den Overbevisning, der fra hans tidlige Ungdom viste
% --- p. 117 ---
% `anti-/rationelle': hyphen at the line end, JOINED: the volume prints
% `antirationel(le)' solid mid-line many times (Striden), never hyphenated;
% witness 2 reads `anti-rationelle'.
\opage{117}sig for ham i Sandhedens Klarhed, han har i de stærkeste anti\-%
rationelle Storme holdt Fornuftens Banner fast.

Hvor mange af Universitetets andre Professorer, f. Eks. af
det naturvidenskabelige Fakultet, har ikke delt hans Synsmaade,
men har, som det er almindelig Skik blandt lærde Folk, for at
undgaa Modstand og Ubehagelighed forholdt sig rolige og altid
forfægtet, at der ingen Grund var til at tale, ‹saa længe man
ikke gjorde dem Noget, ikke gjorde Indgreb i deres Frihed›.
Som om en Mand var en Videnskabsmand, hvis Sandheden ikke
brændte ham som en Glød paa Tungen, som om ikke den, der
tav, lod den Enkelte, der taler, i Stikken og bidrog Sit til at
fremkalde det falske Skin, at han med sin Synsmaade bogstave\-%
% `Livsanskuelse', `urimelig': worn e's that print like c (witness 1
% reads `Livsanskuelsc'); read as e, as throughout the volume.
lig staar ene, at hans Livsanskuelse er en Art Grille, en urime\-%
lig Særhed, opstaaet i en ved falsk Metafysik forskruet Hjerne
osv. Som sagt, Professor Brøchner har ingen Bistand fundet.
Men saa Meget er vist, at er der blandt de Aandsvidenskaberne
studerende Yngre Nogen, som paa sin Universitetsbane har
faaet et stærkt Indtryk af den frie Forsknings ubetingede Ret,
af Humanitetens Gyldighed overfor alle Dogmer, saa skylder
han visselig Professor Brøchner dette Indtryk, og Anmelderen
kan i det mindste vidne for sig selv, at næppe Noget under de
hinanden krydsende Meningers og Systemers forvirrende Kamp
har været ham en saadan Støtte og en saadan Anvisning til at
søge Sandheden ad Fornuftens Vej som Professor Brøchners
strengt videnskabelige Forsken og mandigt ubøjelige Holdning.

Afhandlingen om \emph{Filosofiens historiske Udvikling} er sikkert
det betydeligste af de Skrifter, som Forfatteren endnu har ud\-%
givet. Det er et Værk, som vil blive staaende i vor Literatur.
Det vil være en ypperlig Vejledning for Enhver, som tager fat
paa Studiet af Filosofiens Historie, og det er saa nyt og origi\-%
nalt, saa vel gennemtænkt, atter og atter prøvet, at det endog i
den rige tyske Literatur vilde gøre sig højlig bemærket --- hvil\-%
ken Rolle spiller dets Udgivelse da ikke i en saa fattig filosofisk
Literatur som den danske!

Professor Brøchner har ikke villet give Filosofiens Historie
i og for sig. Han har sat sig et større Maal, han har villet give
denne Histories Filosofi. Hvad vil dette sige? Naar faar Histo\-%
rien sin Filosofi? Den faar det, naar det fyldestgørende paavises,
at den ikke er en Rækkefølge af tilfældige, usammenhængende,
% --- p. 118 ---
\opage{118}lovløse Begivenheder, men naar den ses som lovbestemt. At
den maa være en Lovmæssighed underkastet, godtgør Brøchner
% `dens Ophavs, til Menneskeslægtens Væsen': as printed (both
% witnesses agree); apposition, not a defect.
ved at henvise til dens Ophavs, til Menneskeslægtens Væsen,
hvis Udfoldelse ikke kan være tilfældig eller planløs, men
maa være en virkelig Udvikling. I dette Ord \emph{Udvikling}, det Be\-%
greb, der er Midtpunktet for hele den nyere tyske og engelske
Filosofi, ligger alle de Bestemmelser eller Love sammenrullede,
som en tænkende Betragtning af Historien formaar at udfolde. Har
Historien en Sammenhæng, saa kan ingensomhelst historisk Til\-%
stand, af hvad Art nævnes kan, være uden Sideindflydelse i alle
Retninger; har Historien en Udvikling, saa kan ingen historisk
Kendsgerning, ligegyldig af hvad Art, være uden Eftervirkning
ud i Fremtiden. Tager man et enkelt lille Omraade ud af
Historien, Literaturhistorien f. Eks., da staar denne i Sammen\-%
hæng med alle de øvrige Grene; tager man ud af den igen et
enkelt Afsnit som Filosofiens Historie, da isolerer man paany
kunstigt et enkelt Parti af Historien, der kun forholdsvis udgør et
Hele for sig. Men den Udvikling, som foregaar i det store
Kulturliv, findes ogsaa i dette dets særskilte Organ, ligesom om\-%
vendt den Lovmæssighed, vi finder her, er os en Borgen for, at
der ikke paa noget Punkt kan herske Vilkaarlighed.

De Enkelte, af hvilke Slægten bestaar, mener vel, at de
tænker, som deres egen individuelle Fornuft indgiver dem det,
som de kan og som de vil. I Virkeligheden tænker de dog til\-%
lige og først og fremmest, som de \emph{maa}. De er afledede Væse\-%
ner; de har saa lidt deres Tanker som deres Liv af sig selv;
ligesom det Almene, ligesom Slægten \emph{lever} i dem, saaledes \emph{tæn\-%
ker} ogsaa Menneskeslægten i de Enkelte. De Ideer, den forbe\-%
reder, fremtræder for det første i en Fostertilstand hos de en\-%
kelte Tænkere, men uden at de selv er sig det bevidste, en rum
Tid før disse Ideer, fuldt udviklede, indtager deres Plads i Vi\-%
denskaben. For det andet er Slægtens Udviklingstrin paa det
bestemte Tidspunkt samtidigt som en Skæbne over de Enkelte.
Skønt selve de Ideer, hvis Organer de Enkelte er, paa mange
Maader opfordrer dem til at drage videre Følgeslutninger af
hine Tanker, til at forfølge dem ud i nye, fremtidige Retninger,
er Tidsalderens Aand dog en Magt, der omspænder alle den En\-%
keltes Bestræbelser, ubevidst for ham holder ham tilbage og
forhindrer ham fra at forfølge en Tanke ud over sin Tidsalders
Synskres.

% --- p. 119 ---
\opage{119}Dog den Skæbne, der behersker den Enkeltes Tænkning,
er endnu mere sammensat: ikke blot er han ude af Stand til
at overskride det samtidige Slægtstadiums Grænser, men saa
længe dette Stadium endnu ikke fuldt er modnet og kommet i
Ligevægt om sit eget Midtpunkt, virker det foregaaende, det for\-%
ladte Stadium endnu bestandig som en Magt, der hæmmer den
Enkeltes Stræben. Ogsaa dette er en saavel i Udviklingens Be\-%
greb begrundet som ved historisk Erfaring stadfæstet Lov. Læg\-%
ger man saa sluttelig Mærke til, at tidt indenfor en og samme
historiske Tidsalder stik modsatte Retninger gør sig gældende,
da formaar Historiens Filosof selv her at vise os, hvorledes
det Udviklingsstandpunkt, paa hvilket den hele Tidsalder
staar, uden at de to modsatte Lejre er sig det bevidste, er
fælles for dem begge og forener dem, der anser sig for prin\-%
cipielle Modstandere, under den Alt beherskende Udviklingsform,
der mærkende Ven og Fjende med sit Stempel udformer sig i to
modsatte Skoler som i to Halvkugler, der fuldstændiggør hin\-%
anden.

Det er Love af denne Natur, som Professor Brøchner har
paavist i Filosofiens Historie, og for yderligere at stadfæste sin
Forsknings Resultater har han efter sin Sædvane brugt den
Fremgangsmaade baade at levere Prøve og Modprøve. Han
nøjes ikke med at oplyse sin Tankegang ved Eksempler; men
han vælger, efter at dette er gjort, saadanne Tilfælde, som til\-%
syneladende udgør Undtagelser fra Loven, og viser, hvorledes
den ogsaa finder sin Stadfæstelse i dem.

Hans nye Bogs Fortrinlighed maa dog ikke ene søges i
Lovpaavisningen, ikke heller i Studiet af Forholdet mellem det
i Menneskenaturen Bevidste og Ubevidste, et Omraade, der iøv\-%
rigt noget nær kan betegnes som Professor Brøchners Særom\-%
raade; den beror desuden paa den udtømmende Kundskab til
Filosofiens Historie, ved hvilken Forfatteren indtager saa høj en
Rang mellem Europas tænkende Hoveder. Det er den, som har
gjort ham det muligt at levere saa mesterlige Undersøgelser
som dem af Aristoteles og Kant, der her er forstaaede i en
Dybde og med en Klarhed, som næppe ellers noget Steds. Saa
klar en Tankegang avler en klar og smuk Fremstilling, og der
er derfor store Partier i dette Værk, hvor det er lykkedes For\-%
fatteren trods Æmnets genstridige Natur at bøje det i harmoni\-%
ske Former, og trods sin Pens genstridige Natur nu og da at
% --- p. 120 ---
% `Farve, Saaledes': comma before a capital that opens a new sentence;
% printer's error (a period meant), transcribed as printed.
\opage{120}lade Stemningen meddele Foredraget en mild, aandiggjort Farve,
Saaledes f. Eks. i Skildringen af Nyplatonismen og af Giordano
Bruno.

Til at drøfte Bogens Enkeltheder er her ikke Stedet, selv
om Anmelderen var Opgaven voksen, hvad han ikke er. Der\-%
imod vil jeg rette nogle Indvendinger mod Bogens Plan og imod
Forfatterens Standpunkt.

Planen er ikke aldeles ligelig, og Forfatterens Synsmaade
er ikke ligeoverfor alle Æmner den samme. Han har natur\-%
ligvis Ret til at fremstille Filosofiens Længdesammenhæng for
sig og til at se den som relativt Hele. Men det er tillige klart,
at en saadan Opfattelse med Nødvendighed maa blive meget ab\-%
strakt og temmelig ensidig. Hvad der forekommer mig at være
noget holdningsløst, er dette, at Forfatteren undertiden tager
Bredesammenhængen med og undertiden ikke gør det. Efter
min Opfattelse kan ingen Filosofi fuldt forstaaes, uden at den
ses i dens nære Sammenhæng med alle Literaturens andre
Grene paa dens Tid, med de samtidige historiske Begivenheder
% `Folkeaand, Naar': comma before a new sentence again; printer's
% error, transcribed as printed.
og med Nationens Folkeaand, Naar Talen nu er om den græ\-%
ske Filosofi, saa er Forfatteren villig til at se den i dens Sam\-%
menhæng med den hele græske Aand og Kultur og til at mene,
den kun paa denne Maade kan forstaaes. Men er Talen om
moderne fransk, tysk, engelsk Filosofi, da behandles den, som
om den førte en afsondret Tilværelse uden Forhold til Folkeaan\-%
den, uden Forhold til Landets Poesi, Historie osv. Er det ikke
underligt, slet ikke at indlede Fremstillingen af den tyske Filo\-%
sofi med en Skitse af den tyske Aands Begavelse og Skranker,
men uden videre at betragte den som altomspændende? Den
forekommer maaske Brøchner saaledes, fordi han staar den
saa nær, men er da pur Indskrænkethed Skyld i, at den
aldrig har slaaet an i England? Jeg tvivler i allerhøjeste Grad
derom. Jeg for min Del mener ikke blot, at den i sin Grund
er uforstaaelig uden et foreløbigt Studium af den tyske National\-%
karakter gennem Poesien, Kritiken, Literaturen i det Hele, men
ogsaa at adskillige Særmærker hos dens enkelte Heroer, som
f. Eks. den stærke Hævden af Jeg'et, den stærke Frihedsbegej\-%
string hos Fichte, først ses i deres rette Lys, naar man viser
Napoleons Voldsherredømme over Tyskland i Baggrunden.

Jeg beundrer Brøchners Gave til at fremstille Hovedpunk\-%
terne og Hovedtrækkene i et System; men er hos alle Filosofer
% --- p. 121 ---
\opage{121}Systematiken det Væsenligste og Vigtigste? Er Skelettet Alt, er
Legemet Intet? Bør man ikke i det mindste frygte at stemple
interessante uensartede filosofiske Foreteelser med samme offi\-%
cielle Skoleord, naar rent forskellige Nationalaander er komne
% `fremstilies' for `fremstilles': printer's error (i for l), transcribed
% as printed; clear at 600 dpi (witness 1 agrees, witness 2 corrects).
til Orde i dem? Condillac fremstilies gerne som et Tillæg til
Locke; er der ikke ligefuldt et gabende Svælg mellem deres hele
Maade at filosofere paa? Er der ikke paa Nationalitetens Vegne
et endnu stærkere aandeligt Slægtskab mellem Descartes og Con\-%
dillac, end der med Hensyn til Tænkningens Resultater er mel\-%
lem Condillac og Locke? Ligger Nøglen til Descartes's Filosofi
ikke i hele det store franske Aarhundredes Aand, og hænger
den ikke sammen med Datidens Teologi, Poesi, Lovgivning og
Havekunst? Hører med ét Ord ikke ogsaa Stamme- og Folke\-%
aanden til den Skæbne, som er mægtigere end de Enkeltes, og
som ubevidst indvirker paa dem og hæmmer dem? Det fore\-%
kommer mig, at denne nye Skæbne havde fortjent en stærkere
Betoning.

Og da vi taler om Skæbne, mon da ikke Brøchner selv
har sin? Han tilhører den store Gruppe af Filosofer, der efter
Hegels Tid har forsøgt, med Bibeholdelse af alt det af den store
Mester Indvundne, at bøde paa hans Fejlgreb, undgaa hans
Overgreb og fordybe sig i den Virkelighed, han med sin Ringe\-%
agt for Naturen ofte forsmaaede den erfaringsmæssige Kundskab
til. Saadanne Punkter som Lovundersøgelsen og Paavisningen
af Forholdet mellem det Bevidste og det Ubevidste angiver,
hvor fuldstændigt Brøchner staar paa sin Tidsalders Højde
og deltager i dens Bestræbelser. Imidlertid staar han med Alt
dette jo dog paa Hegels Grund, og jeg vil ikke negte, at jeg er
tilbøjelig til den Tro, at en bevidst-ubevidst Eftervirkning af den
Panlogisme, han har forladt, undertiden holder ham tilbage og
undertiden forleder ham til Forhaands-Konstruktioner, der neppe
holder Stik. Til saadanne regner jeg hans Inddeling af Filoso\-%
fien i tre Hoved-Tidehverv, af hvilke de to er forløbne og den
tredie, Fremtidens Filosofi, nu skal til at begynde. Den Op\-%
gave, som han stiller denne Fremtidens Filosofi, har allerede
det foregaaende Tidehvervs Tænkere stræbt at løse, og om de
i de nærmest følgende Aarhundreder opstaaende Tænkere vil
være heldigere end det indeværendes i at undgaa ensidigt at
lægge Eftertrykket paa Naturen eller paa Aanden, er et Spørgs%
% --- p. 122 ---
\opage{122}maal. Om de endelig vil benytte den Forskningsmetode, Brøch\-%
ner anviser dem, er et endnu større.

Den dialektiske Metode, der hidtil kun er bleven praktiseret
i et enkelt Land og i et kort Tidsrum, har i sin oprindelige
Form ikke staaet Prøve. Den er forladt endog i Tyskland, og
den benyttes ikke i noget andet Land. At det Dialektiske er en
Magt i Menneskeaanden og i Altet, er jeg langt fra at betvivle;
men med hvilken Tillempelse Metoden kan bruges, har Brøch\-%
ner ikke paavist. Alle vi, som arbejder, ønskede hellere end
gerne et fortrinligt Værktøj; men da det højst sindrige, man
hidtil har givet os i Hænde, ikke lader sig bruge, saa kan In\-%
gen tage os det ilde op, at vi foreløbig hellere benytter Fag\-%
videnskabernes grovere Redskaber, som i det mindste ikke brister
i Haanden.

Jeg ser endelig en Levning af Hegels Indflydelse i den
Maade, paa hvilken Brøchner affærdiger den østerlandske Fi\-%
losofi og bestemmer den græske. Den østerlandske regner
han ikke med, en Smule paa Forhaand forekommer det mig.
Der foreligger yderst sammensatte, overordenlig interessante in\-%
diske Systemer, som endnu slet ikke er oversatte; de er, hvis
Anmelderen tør tro Andre paa deres Ord, i deres Væsen filoso\-%
fiske; de har intet særligt Forhold til Religionen, og de maa
dog prøves, førend de kan opofres for den Form af Filosofi,
der som fremstillende ‹den umiddelbare Enhed af Natur og
Aand› nødvendigvis skal være den første.

Til den sidste Bestemmelse kommer Brøchner ad den
Vej, at han begynder med at stemple Grækerne som \emph{Skønhedens
Folk}. Atter dette forekommer mig en utvivlsom Eftervirkning af
Hegelianismen. Saadanne Sætninger er altfor sande, det vil sige,
altfor abstrakte til at være sande. Grækerne er ikke mere Skøn\-%
hedens Folk, end Jøderne er Religionens Folk, end \emph{Don Juan}
er den fuldendte Opera osv. Grækerne udtrykker et enkelt be\-%
stemt historisk Skønhedsideal, og de kom ad rent historisk Vej
til at udtrykke det. Dette Ideal er kun i enkelte Retninger det
højeste indtil nu opstaaede, og det blev dette af rent historiske
Grunde. Grækerne er Opfindere og Stiftere af Byer, der tillige
er Stater. Disse Byers farefulde Beliggenhed blandt kun halvt
undertvungne Stammer, og deres indbyrdes Fejdeliv nødvendig\-%
gjorde den gymnastiske Opdragelse, der atter bevirkede, at den
fuldendte Gymnastiker blev vurderet og beundret i en uhørt Grad,
% --- p. 123 ---
\opage{123}blev bekranset ved offenlige Festligheder, fik Ret til en Statue og
blev Skønhedsidealet. Er da mon ikke den Bestemmelse, at Græ\-%
kerne er Stad-Statens Folk, dybere liggende og mindre ubetinget
end den, at de er Skønhedens? Det forekommer mig, at den er
det, og det synes mig rimeligt, at man ud fra Grækernes gan\-%
ske ejendommelige politiske Liv vilde kunne naa til mange af de
Bestemmelser, der er de mest afgørende for deres Filosofi.

Det er muligt, at jeg har Uret i nogle af disse Indsigelser.
Hvis de imidlertid ikke er ubegrundede, da udelukker de Mang\-%
ler, som de formentlig rammer, lige fuldt paa ingen Maade Bo\-%
gens store Ypperlighed. De godtgør i Grunden kun Gyldigheden
af en af Forfatterens egne Love, nemlig den, at det ældre Tide\-%
hverv, saa længe det nye endnu ikke er fremtraadt i sin fulde
Udvikling, ubevidst for den enkelte Tænker til en vis Grad vir\-%
ker hæmmende ind paa hans bevidste Stræben efter at frem\-%
bringe den nye Tidsalders Filosofi.
\piecerule

\end{document}
